\documentclass[10pt,a4paper]{article}

% ---------- PACKAGES ----------
\usepackage[margin=1.5cm]{geometry}
\usepackage{amsmath, amssymb}
\usepackage{physics}
\usepackage{enumitem}
\usepackage{multicol}
\usepackage{titlesec}
\usepackage{mathrsfs}

% ---------- MISE EN FORME ----------
\setlist[itemize]{noitemsep, topsep=0pt}

\author{Raphaël Jamann}
\title{Formulaire de mécanique - FIMI 2}
\date{}


% ---------- DOCUMENT ----------
\begin{document}
    \maketitle


    \begin{multicols}{2}

    % ======================================================
    \section{Actions mécaniques}

        \subsection*{Moment d'une force}
        \[
        \vec M_F(O) = \vec{OB} \wedge \vec F
        \]

        \subsection*{Transport du moment (BABAR)}
        \[
        \vec M(B) = \vec M(A) + \vec{BA} \wedge \vec R
        \]

        \subsection*{Propriété}
        \[
        \vec M \perp \vec R \quad \text{(glisseur, automoment nul)}
        \]

    % ======================================================
    \section{Torseur d'action mécanique}

        \[
        \{T\}_A =
        \begin{Bmatrix}
        \vec R \\
        \vec M(A)
        \end{Bmatrix}
        \]

        \subsection*{Invariants}
        \begin{itemize}
            \item Invariant vectoriel : $\vec R$
            \item Automoment : $\vec R \cdot \vec M(A) = \text{cte}$
        \end{itemize}

        \subsection*{Torseurs particuliers}
        \begin{itemize}
            \item Glisseur :
            \[
            \begin{Bmatrix}
            \vec R \\
            \vec 0
            \end{Bmatrix}
            \]
            \item Couple :
            \[
            \begin{Bmatrix}
            \vec 0 \\
            \vec C
            \end{Bmatrix}
            \]
        \end{itemize}


    % ======================================================
    \section{Axe central}

        \subsection*{Définition}
        \[
        \vec M(I) = \lambda \vec R
        \]

        \subsection*{Formule}
        \[
        \vec{AI} =
        \frac{\vec R \wedge \vec M(A)}{\|\vec R\|^2}
        + k \vec R
        \]

        \subsection*{Propriétés}
        \begin{itemize}
            \item Moment constant sur l'axe
            \item Moment minimum sur l'axe
        \end{itemize}

    % ======================================================
    \section{Statique}

        \subsection*{Principe Fondamental de la Statique}
        \[
        \sum \vec F_{ext \rightarrow S} = \vec 0
        \qquad
        \sum \vec M_{ext \rightarrow S}(A) = \vec 0
        \]

        \subsection*{Cas plan}
        \[
        \sum F_x = 0 \qquad
        \sum F_y = 0 \qquad
        \sum M_z = 0
        \]

        \subsection*{Cas classiques}
        \begin{itemize}
            \item Solide soumis à 2 forces :
            supports colinéaires, normes égales, sens opposés
            \item Solide soumis à 3 forces coplanaires :
            forces concourantes et somme nulle
        \end{itemize}


    % ======================================================
    \section{Cinématique}

        \subsection*{Repérage d'un point}
        \[
        \vec{OP} = x\vec x + y\vec y + z\vec z
        \]

        \subsection*{Cinématique du point}
        \[
        \vec V_P = \frac{d\vec{OP}}{dt}
        \qquad
        \vec A_P = \frac{d\vec V_P}{dt}
        \]

    % ======================================================
    \section{Cinématique du solide}

        \subsection*{Vitesse}
        \[
        \vec V_P = \vec V_O + \vec\omega \wedge \vec{OP}
        \]

        \subsection*{Accélération}
        \[
        \vec A_P =
        \vec A_O
        + \vec\alpha \wedge \vec{OP}
        + \vec\omega \wedge (\vec\omega \wedge \vec{OP})
        \]

    % ======================================================
    \section{Composition des vitesses}

        \[
        \vec V(P, S/R_0)
        =
        \vec V(P, S/R_1)
        +
        \vec V(P, R_1/R_0)
        \]

    % ======================================================
    \section{Composition des accélérations}

        \[
        \vec A(P, S/R_0)
        =
        \vec A(P, S/R_1)
        +
        \vec A(P, R_1/R_0)
        +
        2 \vec\omega(R_1/R_0) \wedge \vec V(P, S/R_1)
        \]

        \textit{(Terme de Coriolis)}

    % ======================================================
    \section{Dérivée d'un vecteur tournant}

        \[
        \left( \frac{d\vec u}{dt} \right)_{R_0}
        =
        \left( \frac{d\vec u}{dt} \right)_{R_1}
        +
        \vec\omega(R_1/R_0) \wedge \vec u
        \]

    \end{multicols}

\end{document}
